\documentclass[a4paper]{article}
\usepackage[pdftex]{graphicx}
\usepackage[utf8]{inputenc}
\usepackage{enumerate}
\usepackage{href-ul}
\hypersetup{
	colorlinks=true,
	linkcolor=blue,
	urlcolor=blue}
\usepackage{siunitx}
\sisetup{locale=DE} 
\usepackage{icomma}
\usepackage{amssymb}
\usepackage{geometry}
\geometry{a4paper, top=15mm, left=15mm, right=15mm, bottom=15mm,
	headsep=10mm, footskip=12mm}
%Source: img_0108

% Start the document
\begin{document}
%Titel
{\bf Schulaufgabe aus der Physik }\\
\begin{enumerate}[1.]
%Aufgabe 1

{\bf \item Vervollständige die Tabelle :}

\renewcommand{\arraystretch}{3}
\begin{tabular}{|p{3.5 cm}|p{3.5 cm}|p{5 cm}|}
\hline
Physikalische Größe &  Formel & Einheit \\
\hline
& F= $m \cdot a $ & \\
\hline
Druck & &  \\
\hline
 & $ W = F \cdot s $ & \\
\hline
Potentielle Energie & & \\
\hline
& & 1 Watt (1 W)\\
\hline
\end{tabular}

{\bf \item Es wurde für ein Gas die Masse $m = \SI{2,2}{\kilogram}$ und das Volumen $V=\SI{3,1}{\meter^3}$ gemessen.} Berechne die Dichte dieses Gases.  Warum eignet sich dieses Gas zum Füllen von Luftballons?
 ($\varrho_{Luft}= \SI{1,3}{\kilogram \per \meter^3}$)?


{\bf \item Nenne ein Messgerät für den Druck..} \\
a) ...allgemein \hspace{1 cm} b) ...speziell der Luft

{\bf \item Aufgabe}
\begin{enumerate}[a)]
\item Wovon ist der Luftdruck abhängig?
\item wie lautet die Merkregel für den Schweredruck in Wasser? 
\item Berechne die Kraft, die auf eine Fläche von $\SI{600}{\centi\meter^2}$ eines U-Boot-Fensters wirkt, wenn ein Druck von 8 MPa vorhanden ist. 
\item In welcher Tiefe befindet sich das U-Boot?
\end{enumerate}


\end{enumerate} 

\fbox{
	\begin{minipage}{0.5\textwidth}
		Zur Lösung bitte \href{https://www.okuyakl.de/phys/p8druL108/ll108.pdf}{hier klicken} oder den QR-Code scannen.\\
	Weitere Arbeitsblätter gibt es unter 
	
	\href{https://www.okuyakl.de}{www.okuyakl.de}
	\end{minipage}
	\hfill
	\begin{minipage}{0.4\textwidth}
		\includegraphics[width=1.5 cm]{../../viecher/zwe03}
		\includegraphics[width=3 cm]{qr108}
		\includegraphics[width=2 cm]{../../viecher/afanticon1}
		
	\end{minipage}}

\end{document}%Lösung------------------------------------------
